% 第七部：程序面的理由（第伍部分，1386–1395）。出處：英文底稿；W06_1019_導讀問答.md 理論表的調和論證一列（已查核）。
\chapter{程序面的理由：撇開結果，誰有權作終局決定}

\qaentry{第七部之一}{正當性理論須答覆的兩個問題}

\begin{qaquestion}
作者用虛擬的公民 C\textsubscript{n} 提出哪兩個問題？為何正當性理論必須答覆兩者？
\end{qaquestion}

政治決定通常由一組指定的人 \{C\textsubscript{1}, C\textsubscript{2}, \ldots, C\textsubscript{m}\} 以指定的程序作成。C\textsubscript{n} 不在其內，可以問：「Why them? Why not me?」即使接受了答案，C\textsubscript{n} 還可以問：在所用的程序中，為何不給與我看法相同的決定者更大的權重？正當性理論須同時說明成員的組成與所用的程序何以正當\pg{1387}。

\begin{qafigure}{C\textsubscript{n} 的兩個問題與兩個機關的答覆（1387--1392）。\label{fig:cn}}
% C_n 的兩個問題 × 兩個機關（1387–1392）。三格灰階，第二問對法院一格用強調色。
\begin{tikzpicture}[qa,cell/.style={nbox,text width=52mm,minimum height=15mm}]
  \node[hdr] at (3.15,1.2) {立法機關（約五百人）};
  \node[hdr] at (8.85,1.2) {聯邦最高法院（九人）};
  \node[hdr,anchor=east,align=right] at (0.15,0) {第一問\\[-1pt]{\scriptsize\mdseries 為何由這些人}};
  \node[hdr,anchor=east,align=right] at (0.15,-1.8) {第二問\\[-1pt]{\scriptsize\mdseries 為何採此程序}};
  \node[cell] at (3.15,0) {公平選舉：C\textsubscript{n} 與全體公民受到平等對待（1387）};
  \node[cell] at (8.85,0) {總統提名、參議院同意的間接民主資格，與議員相較遠為不足（1391）};
  \node[cell] at (3.15,-1.8) {多數決的公平性論證：中立、平等，給每一意見最大可能的權重（1388）};
  \node[cell,kbox] at (8.85,-1.8) {五票勝過四票；大法官不代表任何人，多數決在此欠缺道德基礎（1391--1392）};
\end{tikzpicture}

\end{qafigure}

\qaentry{第七部之二}{立法機關的答覆}

\begin{qaquestion}
對立法機關，C\textsubscript{n} 的兩個問題如何得到答覆？
\end{qaquestion}

作者把兩個問題寫成具體的形式：為何由大約五百位議員決定一個影響我與另外兩億五千萬人的權利問題？即使接受，為何不給與我看法相同的議員更大的權重\pg{1387}？第一問由公平選舉的理論答覆，C\textsubscript{n} 在決定誰得參與這類決定時，與全體公民受到平等對待；第二問由多數決（MD）的公平性論證答覆\pg{1387--1388}。作者寫道，MD 比其他規則更能在爭議的結果之間保持中立、平等對待參與者，並在給予所有意見同等權重的前提下，給每一意見最大可能的權重\pg{1388}。代表之間的多數決與選舉合起來，近似全體公民之間的多數決，給予平等的發言權與平等的決定權威\pg{1388--1389}。現實制度的缺陷由第一項假設處理，作者寫，正當性理論「will talk about reasonable fairness, not perfect fairness」\pg{1389}。

\qaentry{第七部之三}{法院的答覆與 Scalia 的引文}

\begin{qaquestion}
同樣兩個問題向聯邦最高法院提出時，答覆為何？作者為何長段引用 Scalia？
\end{qaquestion}

作者說正當性是比較性的，於是向法院提出 C\textsubscript{n} 的兩個問題\pg{1389}：為何由這九人決定此事？即使由他們決定，為何採用現行的程序，而不採給予與 C\textsubscript{n} 看法相同的大法官更大權重的程序\pg{1390}？作者寫，這兩問對法院遠比對議員難答，接著長段引用 Scalia 大法官在 Planned Parenthood v.\ Casey 案的不同意見：Scalia 對遊行、信件、抗議等「政治壓力」同感困擾，認為成因在於法院不再依文本與傳統實務，而依本院所稱的「說理的判斷」，那不過是哲學偏好與道德直覺\pg{1390--1391}。

第一問的答覆是大法官經總統提名、參議院同意而有間接的民主資格；作者認為與議員相較遠為不足\pg{1391}。第二問的處境更糟：

\begin{quote}
But in the end it comes down to head-counting: five votes defeat four in the U.S. Supreme Court, irrespective of the arguments that the Justices have concocted. \ldots{} MD is appropriate for persons who have a moral claim to insist on being regarded as equals in some decision-process. But I cannot see any moral basis for this claim in the case of Supreme Court Justices. They do not represent anybody. Their claim to participate is functional, not a matter of entitlement.\pg{1391--1392}
\end{quote}

文獻對法院內部的多數決少有討論；若以 Condorcet 陪審團定理為法院辯護，該辯護屬結果面理由，而依該定理，群體愈大，增進愈多，反而有利於人數較多的立法機關\pg{1392~及註~122}。

\qaentry{第七部之四}{五項調和論證及其反駁}

\begin{qaquestion}
辯護者提出哪些論證，主張違憲審查與民主可以調和？\\
作者如何逐一反駁？
\end{qaquestion}

\begin{longtable}{@{}L{0.42\linewidth}L{0.50\linewidth}@{}}
\toprule
調和論證 & 作者的反駁 \\
\midrule
\endhead
法官只執行有民主資格的權利法案 & 權利法案的抽象條文不能解決分歧\pg{1393} \\
預先承諾：社會像 Ulysses 把自己綁在桅桿上，法官是確保繩索未鬆的船員 & 社會並未承諾任何特定的權利理解；一旦出現另一種理解，既有的預先承諾為何仍應拘束並不清楚\pg{1393--1394}；註 123 引 Elster, Ulysses Unbound 對先前 Ulysses and the Sirens 的修正 \\
得修改權利法案而未修改，等於默示認可 & 修改通常需要特別多數；英國或紐西蘭式的法律位階權利法案也因其在政治文化中的地位而提高修改的負擔；起點由判決設定，不修改不能視為默示認可\pg{1394} \\
法官經提名同意而具有民主資格（Eisgruber） & 法官的民主資格遠不及議員；以民主資格接受法官，反而削弱把違憲審查當作一種獨特而有價值之決策形式的正面論證\pg{1394} \\
訴訟是公民參與的額外途徑 & 這種途徑不受政治平等原則的約束\pg{1395} \\
\bottomrule
\end{longtable}

\begin{readernote}
導讀時可以點出第四項的反駁：辯護者一旦以法官的民主資格為理由，就承認了程序面的比較，而在這項比較上法官必然輸給議員。作者據此說它削弱了違憲審查自己的正面論證。
\end{readernote}
